\section{Results}
\label{sec:results}

% =============================================================================
% SKELETON --- restructured 2026-09-21 to the six-section layout of
% SOL_PLAN_REFORMAT.md. Written ONLY from the final run (20 replicates, shortfall path
% capped, shock probability re-fitted to 0.040 on the rebuilt population). Nothing in
% results/ dated before that run may be quoted: it predates DEFECTS.md B32.
%
% What moved out of this section:
%   - the baseline and RQ1          -> Section~\ref{sec:data-validation} (Calibration)
%   - the three scenarios and RQ3   -> Section~\ref{sec:scenarios}
%
% BUDGET: 2,100 words of PROSE. Tables and figures do not count, so every number
% goes in a table and prose is spent only on what the number means.
%
% Section budgets, which are binding:
%   opening                          100
%   4.1 Invisible credit and stacking 550
%   4.2 Access, peer influence, default 750   (includes the written no-cascade paragraph)
%   4.3 Distribution and mechanism    350
%   4.4 Sensitivity and robustness    350
%
% OPENING (100 words): start from the calibrated benchmark of Section~\ref{sec:lim-calibration};
% state the comparison protocol once --- twenty replicates, the seed scheme, the noise
% floor (report the final-run replicate sd; differences under two sd are noise), and
% that access rates are shares of the BANKED subpopulation (82.8% ceiling) --- then a
% one-sentence roadmap. State it here so no subsection has to repeat it.
%
% EVERY table and figure carries a short \caption and a \tabnote that states: the
% replicate count and seed scheme; the noise floor; which scenario each column is;
% that access rates are shares of the BANKED subpopulation, not of all households; and
% that beta = 0 is the control arm. This is a caption requirement, not a prose habit.
%
% Report a result, contrast it with its control, explain the mechanism. In that order.
% =============================================================================

\subsection{Invisible Credit and Stacking}
\label{sec:results-stacking}

% BUDGET: 550 words. The stacking half of RQ2.
%
% CONTENT REQUIRED:
% - Report the no-BNPL baseline BEFORE any BNPL-on outcome.
% - Stacking is EMERGENT, not programmed. Share of adopters holding 2+ facilities
%   against the CFPB's 32%, as an order-of-magnitude check only.
% - Platform-count sweep 1-6: the single-platform arm gives exactly zero cross-platform
%   stacking, which isolates cross-platform accumulation from single-platform
%   accumulation; then how concurrent facilities and default change as the count rises.
% - This subsection is the EVIDENCE BASE for the argument of Section~\ref{sec:scenarios}:
%   the harm is the number of agreements a household holds. Make the three facts easy
%   to cite from there (2+ share, zero with one platform, default rising with count).
% - Figure: rq1_stacking.png
%   NOTE: figure FILENAMES keep the original RQ numbering and are deliberately frozen
%   (they are baked into results/raw/ labels and analysis.py load keys). rq1_stacking.png
%   belongs to thesis RQ2. See the numbering note at the top of simulation/experiments.py.

\subsection{BNPL Access, Peer Influence and Default}
\label{sec:results-rq2}

% BUDGET: 750 words, of which the paragraph below (about 130) is written.
%
% CONTENT REQUIRED:
% - Default with BNPL off, on, and on with strong peer influence. The beta=0 row must
%   be explicit in every table and on every figure. Separate the FINANCIAL channel
%   from the SOCIAL channel: beta is the strength of peer influence on adoption, so
%   sweeping it is what separates a social explanation from a purely financial one.
% - Access sweep runs inside the banked subpopulation (82.8% ceiling, Submodel~12).
%
% - THE NEGATIVE RESULT, with its own bold lead sentence, not buried:
%   the increase is SMOOTH in access. No cliff edge anywhere on the surface. This
%   falsifies the pre-registered non-linearity hypothesis.
%   Present the linear and the pre-registered Granovetter threshold specifications
%   TOGETHER: the threshold variant was run precisely because a linear rule might give
%   a linear answer for trivial reasons. Report what adoption does and what default
%   does under each. Two structurally distinct social mechanisms.
%   The structural reason is the paragraph below, which is already written and must
%   sit directly after the negative result.
%
% - Pattern 1 is tested here: what enabling BNPL does to traditional servicing stress.
%   Attach the caveat of Section~\ref{sec:lim-validation} --- the direction is partly
%   designed into the lender-choice rule, the magnitude is not.
% - Figure: rq2_default_surface.png
% - Table: default rate by access x beta, with the beta=0 row first.

% --- Written prose: the structural explanation of the negative result. Keep directly
% after the result it explains. ---
\label{sec:no-cascade}
The absence of a cliff edge is a property of a model built this way rather than a general finding.
Six mechanisms are excluded by design (Table~\ref{tab:design-concepts}), and the second of them, the absence of any channel by which one household's default affects another household's income, is the explanation of the negative result.
With no such channel, what the model produces is many correlated individual defaults rather than a cascade, and a threshold in the population default rate requires feedback in the outcome itself: households catch \ac{BNPL} from their neighbours, but they do not catch financial distress from them.
A successor design with an impairable lender or a default-to-income channel could plausibly produce a cliff edge, so the negative result is bounded rather than general.

\subsection{Distribution and Mechanism}
\label{sec:results-distribution}

% BUDGET: 350 words.
%
% CONTENT REQUIRED:
% - Where debt service and default concentrate, by income quintile. The empirical
%   record places BNPL harm among the liquidity-constrained; say whether the model does.
% - THE SPLIT ANSWER, reported honestly and in ONE place: how often the rolling limit
%   binds, and what that does to default. In the working run the limit bound on about
%   half of requests yet barely moved default, because a constrained household borrows
%   less per purchase rather than defaulting less. Reported apart, those two facts read
%   as a contradiction. Re-check both against the final run before writing a word.
% - State which direction is partly induced by the lender-choice rule and which
%   magnitude is genuinely produced by the model.

\subsection{Sensitivity and Robustness}
\label{sec:results-robustness}

% BUDGET: 350 words. One figure, one table, prose only for what moves the answer.
%
% CONTENT REQUIRED:
% - Figure: robustness_tornado.png. Each bar is how far default moves across the
%   plausible range of one parameter; red bars are rules with no citation.
% - The Submodel~4 borrowing-amount rule (Section~\ref{sec:amount}) was the largest
%   sensitivity in the working run, with the whole range coming from the most generous
%   of the three specifications. The shortfall path is now CAPPED (Submodel~5), and the
%   robustness suite runs all three rules capped and uncapped on the same seeds: state
%   plainly whether the cap narrows the sensitivity, and by how much.
% - Confirm empirically, rather than asserting, that the peer channel is insensitive to
%   activation order.
% - Everything else --- Sobol indices, activation-order re-runs, population-size
%   stability at 1,000/5,000/10,000, the minimum-payer sweep 0.20-0.40, and the full
%   platform-count detail --- goes to Appendix~\ref{app:supplementary} as tables, with
%   one sentence here saying so. Do not narrate them.
